\documentclass[10pt,twocolumn]{article}
\usepackage[T1]{fontenc}
\usepackage[utf8]{inputenc}
\usepackage{lmodern}
\usepackage[margin=1.8cm,top=2.4cm,bottom=2cm,columnsep=0.6cm]{geometry}
\usepackage{fancyhdr}
\usepackage{titlesec}
\usepackage{booktabs}
\usepackage{amsmath,amssymb,amsfonts}
\usepackage{microtype}
\usepackage{xcolor}
\usepackage{mdframed}
\usepackage{parskip}
\usepackage{enumitem}
\usepackage{hyperref}

\definecolor{rulered}{RGB}{180,30,30}
\definecolor{boxgray}{RGB}{245,245,245}
\definecolor{keywordgray}{RGB}{100,100,100}

\pagestyle{fancy}
\fancyhf{}
\fancyhead[L]{\textsc{\small Claude Mag}}
\fancyhead[C]{\small\textit{Machine Learning Research Digest}}
\fancyhead[R]{\small 2026-11-11}
\fancyfoot[C]{\small\thepage}
\renewcommand{\headrulewidth}{0.8pt}

\titleformat{\section}{\normalsize\bfseries\color{rulered}}{}{0pt}{}%
  [\vspace{-4pt}\color{rulered}\rule{\columnwidth}{0.4pt}]
\titlespacing{\section}{0pt}{8pt}{4pt}

\hypersetup{colorlinks=true,urlcolor=rulered,linkcolor=black}

\begin{document}
\setlength{\parindent}{0pt}
\setlength{\parskip}{4pt}

{\small\color{keywordgray}\textbf{Keywords:}
  sparse autoencoders \textbullet{}
  mechanistic interpretability \textbullet{}
  scaling laws \textbullet{}
  activation buffer}\par\vspace{4pt}

{\LARGE\bfseries\raggedright
  Scaling Laws for SAE Training Data}\par\vspace{4pt}

{\large\itshape\raggedright
  SAE Quality Depends on Buffer Diversity, Not Fresh Tokens: 8--64$\times$
  Storage Reduction with Negligible Quality Loss}\par\vspace{6pt}

{\small\textbf{By} Nikita Koriagin, Nikita Balagansky, Daniil Gavrilov
  (T-Tech)}\par\vspace{2pt}

{\small\color{keywordgray}NeurIPS 2026 Poster \#148634 \textbullet{}
  also ICML 2026}
\par\vspace{2pt}\noindent{\color{rulered}\rule{\columnwidth}{0.6pt}}\vspace{6pt}

Training large sparse autoencoders is bottlenecked by activation storage.
The authors show that SAE quality depends on the number of unique activations
in the buffer --- not on whether each training step uses fresh tokens. After
a short diversity-limited regime, replay of a fixed buffer is as good as
continuous fresh-data collection, cutting storage requirements by 8--64$\times$
with negligible quality loss across two model families, three layers,
three SAE architectures, and four evaluation metrics.

\begin{mdframed}[backgroundcolor=boxgray,linecolor=rulered,linewidth=1pt,
    innerleftmargin=6pt,innerrightmargin=6pt,
    innertopmargin=6pt,innerbottommargin=6pt]
\textbf{\small AT A GLANCE}\par\vspace{4pt}
\begin{description}[leftmargin=0pt,labelsep=4pt,itemsep=2pt,parsep=0pt,
    font=\small\bfseries]
  \item[Problem:] SAE training requires storing billions of LLM activations;
    the activation buffer dominates I/O cost and disk usage.
  \item[Why it matters:] SAEs are a leading interpretability tool; reducing their
    training cost unlocks larger-scale and more accessible interpretability
    research.
  \item[Method:] Fit a data-constrained scaling law; identify the unique-activation
    threshold $n^*$; collect $n^*$ diverse activations once, then replay.
  \item[Result:] 8--64$\times$ reduction in activation storage; $<1\%$ quality
    loss on reconstruction and interpretability metrics.
  \item[Evidence:] Llama-3.1-8B and Qwen3-7B; early/middle/late layers;
    BatchTopK, TopK, JumpReLU SAEs; 50M--10B token budgets.
\end{description}
\end{mdframed}

\section{The Big Picture}

Sparse autoencoders (SAEs) have become the workhorse of mechanistic
interpretability: trained on the residual stream activations of a frozen LLM,
they decompose those activations into sparse linear combinations of learned
features, each of which often corresponds to an interpretable concept.

Training a high-quality SAE requires millions to billions of activation examples.
Generating these activations requires running the LLM forward on a large text
corpus, which produces enormous data volumes: a single residual stream layer of
Llama-3.1-8B at fp16 precision generates 16 KB per token, so 1B tokens require
16 TB of storage. This activation bottleneck is a major practical obstacle to
large-scale SAE research, particularly for researchers without datacenter-scale
storage infrastructure.

The standard recipe --- generate fresh activations for every training step from
a streaming corpus --- implicitly assumes that data novelty is the scarce
resource. The central finding of this paper is that this assumption is wrong for
SAEs: the scarce resource is \emph{activation diversity} (coverage of the
activation distribution), not novelty (seeing each activation only once).
Once the buffer achieves adequate coverage, additional unique activations provide
negligible benefit, and the buffer can be replayed indefinitely.

\section{Why Existing Approaches Fall Short}

\textbf{Streaming activation generation} (the standard recipe) generates fresh
activations from a streaming corpus for every training step. It never repeats
any activation, maximizing novelty. It is correct but expensive: it ties
training speed to LLM inference throughput and requires either online inference
or offline storage at full token budget.

\textbf{Fixed small buffers without principled sizing} have been used in practice
(e.g., the OpenAI SAE release uses a buffer of 400K activations) but without a
principled criterion for the right buffer size. Too small a buffer leads to
mode collapse; too large is unnecessarily costly.

\textbf{Curriculum-based data selection} has been studied for LLM training but
not for SAEs, whose training objective (sparse reconstruction) has very different
data-efficiency properties from next-token prediction.

This paper is the first to characterize the data-efficiency of SAE training
formally, through the lens of a data-constrained scaling law, and to derive the
optimal buffer-sizing strategy from first principles.

\section{Core Mechanism}

The key observation is that two SAE training runs with equal numbers of
\emph{unique} activations but different levels of repetition (one streaming,
one with replay) produce SAEs of comparable quality. This decouples the training
budget (number of gradient steps) from the data budget (number of unique
activations).

The unique-activation count $n_\text{unique}$ is the number of distinct
activation vectors that the SAE has seen at least once during training. The
authors find empirically that SAE quality $\mathcal{L}$ is well described by:
\[
  \mathcal{L}(n_\text{unique}) = \mathcal{L}_\infty + C \cdot n_\text{unique}^{-\alpha},
\]
where $\mathcal{L}_\infty$ is the asymptotic quality, $C > 0$ is a scaling
constant, and $\alpha > 0$ is the data-scaling exponent. This power law holds
across model families, layer positions, SAE architectures, and evaluation metrics.

\section{Technical Formulation}

SAE training minimizes:
\[
  \mathcal{L}_\text{SAE}(W_e, W_d, b) =
    \mathbb{E}_x \bigl[ \|x - W_d\,\text{TopK}(W_e x + b)\|_2^2
    + \lambda \|\text{TopK}(W_e x + b)\|_1 \bigr],
\]
where $x \in \mathbb{R}^d$ is an activation vector, $W_e \in \mathbb{R}^{m\times d}$
is the encoder, $W_d \in \mathbb{R}^{d\times m}$ is the decoder, and
$\text{TopK}$ selects the $k$ largest pre-activations and zeros the rest.

The \textbf{data-constrained scaling law} is fit by:
\begin{enumerate}[itemsep=2pt,parsep=0pt]
  \item Train multiple SAEs with the same gradient-step budget $T$ but varying
    buffer sizes $n_1 < n_2 < \cdots < n_B$.
  \item Measure final quality $\mathcal{L}_i$ for each run.
  \item Fit the three parameters $(\mathcal{L}_\infty, C, \alpha)$ via
    non-linear least squares.
\end{enumerate}
The optimal buffer size $n^*$ is defined as the smallest $n_\text{unique}$ such
that $\mathcal{L}(n^*) \leq \mathcal{L}_\infty + \delta$ for a chosen tolerance
$\delta$ (e.g., 1\% above $\mathcal{L}_\infty$). The budget-optimal training
strategy then collects $n^*$ unique activations once and replays them uniformly
at random for the remaining $T - n^*$ gradient steps.

\section{Learning or Inference Procedure}

The practical buffer-sizing workflow is:
\begin{enumerate}[itemsep=2pt,parsep=0pt]
  \item Run a small calibration sweep: train 5--10 SAEs with varying buffer sizes
    at 1\% of the full training budget.
  \item Fit the scaling law to the calibration sweep results.
  \item Identify $n^*$ from the fitted law.
  \item Collect $n^*$ activations from the LLM (a one-time cost); store this
    buffer on disk.
  \item Train the full SAE by replaying activations sampled uniformly from the
    buffer.
\end{enumerate}
The calibration sweep adds 10\% overhead relative to a single training run;
this is recovered in full by the reduced storage and I/O cost.

\section{What the Guarantee Says}

No formal convergence or approximation bound is derived. The empirical claim is:
across all evaluated settings (Llama-3.1-8B, Qwen3-7B; layers 4, 16, 28;
BatchTopK/TopK/JumpReLU; reconstruction loss, dead-neuron fraction, feature
absorption, automated interpretability), training with the optimal-buffer-size
replay strategy degrades quality by less than 1\% relative to the streaming
baseline at a storage reduction of 8--64$\times$.

\section{Experimental Findings}

\begin{table}[h!]
\centering\small
\begin{tabular}{lcc}
\toprule
Setting & Storage reduction & Quality loss \\
\midrule
Llama-3.1 layer 4 & 64$\times$ & $<$0.5\% \\
Llama-3.1 layer 16 & 32$\times$ & $<$0.3\% \\
Llama-3.1 layer 28 & 16$\times$ & $<$0.4\% \\
Qwen3 all layers & 8--32$\times$ & $<$1\% \\
Automated interp. & 16$\times$ & $<$2\% \\
\bottomrule
\end{tabular}
\end{table}

The largest storage reductions occur in early layers, where activations are more
stereotyped and coverage is achieved with fewer unique examples. Late layers,
where activations carry more semantic diversity, require more unique examples
but still admit substantial savings. The automated interpretability score is
the most variable metric but still achieves 16$\times$ reduction within tolerance.

\section{Ablations and Interpretation}

\textbf{Scaling law fit quality:} The power-law form fits observed data with
$R^2 > 0.98$ in all 18 evaluated settings (3 layers $\times$ 3 architectures
$\times$ 2 models), confirming the law is not an artifact of a particular setup.

\textbf{Scaling exponent $\alpha$:} Values range from 0.3 (early layers, Llama-3.1)
to 0.7 (late layers, Qwen3). Higher $\alpha$ means quality improves more steeply
with $n_\text{unique}$; these settings benefit less from replay.

\textbf{Uniform vs.\ prioritized replay:} Replaying activations proportional to
their reconstruction loss (prioritizing hard examples) improves quality by
$<$0.2\% over uniform replay at the cost of maintaining priority weights;
uniform replay is recommended for simplicity.

\textbf{Generalization to new domains:} A SAE trained with a Wikipedia-only buffer
and then deployed on code activations degrades by 8--12\%; this confirms the
unique-activation hypothesis is domain-specific and buffer diversity must match
the deployment distribution.

\section*{Reference}

Nikita Koriagin, Nikita Balagansky, Daniil Gavrilov.
\textbf{Scaling Laws for SAE Training Data.}
NeurIPS 2026 Poster \#148634.
\url{https://neurips.cc/virtual/2026/poster/148634}

\end{document}
